\documentclass[11pt,a4paper]{report}
\usepackage[utf8]{inputenc}
\usepackage[T1]{fontenc}
\usepackage{lmodern,microtype}
\usepackage[margin=2.4cm]{geometry}
\usepackage{graphicx}
\usepackage{booktabs,longtable,array,amsmath,amssymb,xcolor,listings,hyperref,enumitem,titlesec,etoolbox}
\hypersetup{colorlinks=true,linkcolor=blue!50!black,urlcolor=blue!50!black,citecolor=blue!50!black,pdftitle={nwdynam: Dynamic network actor models in Stata. Manual},pdfauthor={Thomas U. Grund}}
\lstdefinestyle{stata}{basicstyle=\ttfamily\footnotesize,backgroundcolor=\color{black!4},frame=single,rulecolor=\color{black!20},breaklines=true,columns=fullflexible,keepspaces=true,xleftmargin=2pt,xrightmargin=2pt,aboveskip=5pt,belowskip=5pt}
\lstset{style=stata}
\setlength{\emergencystretch}{3em}
\newcommand{\nwdynam}{\texttt{nwdynam}}
\newcommand{\goldfish}{\texttt{goldfish}}
\newcommand{\gfver}{1.6.12}
% \opt sets an option name in typewriter type and lets long names break after ( : , and _ so that they do not run into the next table column
\ExplSyntaxOn
\tl_new:N \l__nwdynam_opt_tl
\NewDocumentCommand{\opt}{m}
  {
    \tl_set:Nn \l__nwdynam_opt_tl {#1}
    \regex_replace_all:nnN { ([\:\,]) } { \1 \c{allowbreak}\cB\{\cE\} } \l__nwdynam_opt_tl
    \regex_replace_all:nnN { \c{\_} } { \c{\_} \c{allowbreak}\cB\{\cE\} } \l__nwdynam_opt_tl
    \regex_replace_all:nnN { (/) } { \1 \c{allowbreak}\cB\{\cE\} } \l__nwdynam_opt_tl
    \regex_replace_all:nnN { ([\|\(\[=]) } { \1 \c{allowbreak}\cB\{\cE\} } \l__nwdynam_opt_tl
    \texttt{ \tl_use:N \l__nwdynam_opt_tl }
  }
\ExplSyntaxOff
% the listing of the output of an example (manual/ex/out/<name>.txt, cut from the Stata logs by ex/mklistings.py)
\newcommand{\showlog}[1]{\lstinputlisting{ex/out/#1.txt}}
\newcommand{\showlogpart}[3]{\lstinputlisting[firstline=#2,lastline=#3]{ex/out/#1.txt}}
\newcommand{\inc}[1]{\typeout{>>>#1}\input{#1}}
\newcommand{\cmpMaxB}{5.6e-07}
\newcommand{\cmpMaxSE}{4.6e-08}
\newcommand{\cmpMaxLL}{5.5e-14}
\newcommand{\cmpNModels}{5}
\newcommand{\cmpMaxWhere}{Emails, choice, \texttt{trans}}

\newcommand{\manualseries}{1.11}
% The manual is labelled by the minor series (valid for 1.11.x): change \manualseries only for a new minor or major release whose documented behaviour changed.
\newcommand{\pkgversion}{1.11.2}
\newcommand{\pkgdate}{26 November 2026}
% Exact version and date of the release that produced the listings (tools/set_version.py does not touch this file; change by hand and rerun the examples).
\newcommand{\manualurl}{https://nwcommands.org/nwdynam/nwdynam_manual.pdf}
\newcommand{\manualdoi}{10.5281/zenodo.24398099}
\newcommand{\citedoi}{ Zenodo. \url{https://doi.org/\manualdoi}}
\title{\textbf{nwdynam}\\[4pt]\large Dynamic network actor models in Stata\\ (a package that builds on \texttt{nwcommands})\\[6pt]\normalsize User manual and reference}
\author{Thomas U. Grund\\ RWTH Aachen University}
\date{This manual describes \texttt{nwdynam} \manualseries.x\\ (examples produced with \texttt{nwdynam} \pkgversion, \pkgdate)}
\begin{document}
\maketitle
\begin{quote}\noindent\textbf{How to cite.} This manual describes the series~\manualseries{} of \nwdynam{} (it holds for all versions \manualseries.x); the listings were produced with version~\pkgversion{} (\pkgdate). It is the citable reference for the package. Please cite it as:\par\smallskip
\noindent Grund, T.~U. (2026). \emph{nwdynam: Dynamic network actor models in Stata. Manual.}\citedoi\par\smallskip
\noindent Please also report the exact version of the program that produced your results, as printed by \texttt{nwdynam version} and stored in \texttt{e(version)}.
\end{quote}
\chapter*{Acknowledgments}\addcontentsline{toc}{chapter}{Acknowledgments}
\nwdynam{} follows the \goldfish{} project closely and exists because of it. The dynamic network actor model is due to Christoph Stadtfeld and Per Block~\cite{stadtfeld2017b}, the coordination model to Stadtfeld, Hollway and Block~\cite{stadtfeld2017h}, and \goldfish{} (James Hollway, Christoph Stadtfeld, Marion Hoffman, Alvaro Uzaheta and contributors~\cite{goldfish}) is the reference implementation on which the effects, the likelihoods, the defaults and the postestimation tools of \nwdynam{} are modeled. The R source code and the documentation of \goldfish{}~\gfver{} were the specification against which the Stata code was written and checked. The change point methods of \texttt{estat changepoints} are a port of the R package \texttt{changepoint} (Rebecca Killick, Kaylea Haynes, Idris Eckley, Paul Fearnhead and Jamie Lee~\cite{killick2014}), which \goldfish{} calls for the same purpose.

\nwdynam{} is not affiliated with or endorsed by the \goldfish{} project, and its errors are the author's. It is meant for researchers who work in Stata or prefer it to R. It is not a replacement for \goldfish{}, which remains the reference.

\paragraph{Use of an AI coding assistant.} The software was developed with Claude Code (Anthropic), an AI coding assistant, working under the author's direction. It helped port algorithms from \goldfish{} into Stata, Mata and C and helped write the tests, and it helped draft this manual. The author specified the scope, reviewed the code, and is responsible for the package and for the manual. Its correctness rests on validation, not on the assistant: the results are compared numerically with \goldfish{}, and the test suite reruns these comparisons.

\paragraph{License.} \nwdynam{} is free software under the GNU General Public License, version~3; Thomas U.\ Grund is its sole author.

\tableofcontents
\inc{ch01_intro}
\inc{ch02_model}
\inc{ch03_start}
\inc{ch04_data}
\inc{ch05_options}
\inc{ch06_effects}
\inc{ch07_netterms}
\inc{ch08_special}
\inc{ch09_postest}
\inc{ch10_examples}
\inc{ch11_validation}
\inc{ch12_limits}
\inc{ch13_technical}
\appendix
\inc{app_translation}
\bibliographystyle{plain}
\begin{thebibliography}{99}
\bibitem{stadtfeld2017b} Stadtfeld, C., Block, P. (2017). Interactions, actors, and time: Dynamic network actor models for relational events. \emph{Sociological Science} 4, 318--352. \url{https://doi.org/10.15195/v4.a14}
\bibitem{stadtfeld2017h} Stadtfeld, C., Hollway, J., Block, P. (2017). Dynamic network actor models: Investigating coordination ties through time. \emph{Sociological Methodology} 47(1), 1--40.
\bibitem{goldfish} Hollway, J., Stadtfeld, C., Hoffman, M., Uzaheta, A. and contributors. \emph{goldfish: Statistical Network Models for Dynamic Network Data}. R package. \url{https://cran.r-project.org/package=goldfish} (compared with version \gfver).
\bibitem{killick2014} Killick, R., Eckley, I.A. (2014). changepoint: An R package for changepoint analysis. \emph{Journal of Statistical Software} 58(3), 1--19.
\bibitem{butts2008} Butts, C.T. (2008). A relational event framework for social action. \emph{Sociological Methodology} 38, 155--200.
\bibitem{michalski2020} Michalski, R. (2020). Manufacturing company email metadata and corporate hierarchy. Harvard Dataverse, V2. \url{https://doi.org/10.7910/DVN/6Z3CGX}. Released under CC0 1.0.
\bibitem{nurek2020} Nurek, M., Michalski, R. (2020). Combining machine learning and social network analysis to reveal the organizational structures. \emph{Applied Sciences} 10(5), 1699.
\bibitem{madan2012} Madan, A., Cebrian, M., Moturu, S., Farrahi, K., Pentland, A. (2012). Sensing the ``health state'' of a community. \emph{IEEE Pervasive Computing} 11(4), 36--45. % TODO(owner): confirm that this is the reference goldfish gives for its Social_Evolution data
\end{thebibliography}
\end{document}
